\section{Open Questions (5)}

\begin{enumerate}

\item \textbf{Block-realisation dependence of the $3/8$ decay base.} Our MC-fitted base $0.282$ is $\sim 25\%$ below the paper's Thm~3 prediction $3/8 = 0.375$, using an $RX{\cdot}RZ{\cdot}CX{\cdot}RX{\cdot}RZ$ 8-parameter 2-qubit block instead of the paper's literal Eq.~10 block. Both are locally 2-design-equivalent at each stage, but the numeric prefactor is evidently block-realisation-dependent. \emph{Next steps:} symbolically re-derive Eq.~13 for the 8-parameter block via Haar integration over its parameter manifold; run MC with three block families (paper Eq.~10, our block, an $RY{\cdot}CZ{\cdot}RY$ HEA-style block) on the same $N=3\ldots 9$ grid and compare fitted bases; predict analytically whether the base depends on 2-design order or on higher Haar moments.

\item \textbf{qTTN / qMERA polynomial-scaling reproduction (C4).} The paper's most consequential structural claim --- that tree/MERA-inspired ansatze exchange the qMPS exponential decay for a polynomial one --- was \emph{not} exercised here. A REPLICATED verdict on C1+C2 alone does not vouch for the polynomial-vs-exponential dichotomy. \emph{Next steps:} implement qTTN and qMERA staircase ansatze in PennyLane (tree structure with log-depth for TTN, MERA disentanglers + isometries per layer); run the MC estimator on $N$ up to 16 (feasible on \texttt{default.qubit} for TTN, marginal for MERA); fit $\log\mathrm{Var}$ vs $\log N$; extract the polynomial exponent; compare to Thm 4/5 predictions.

\item \textbf{Interaction with problem-inspired / adaptive initialisation.} The paper analyses uniform-Haar-random initialisation $\theta \sim \mathrm{Unif}[-\pi,\pi]$. Real VQE/QAOA workflows initialise closer to identity or via layerwise pretraining, and the plateau's dependence on non-uniform initialisation distributions is not addressed. \emph{Next steps:} rerun the MC estimator with (a) identity-close init ($\theta \sim \mathcal{N}(0,\sigma)$ small $\sigma$), (b) beta-distribution init concentrated near 0, (c) layerwise-pretrained init from a small VQE run on a toy Hamiltonian; plot Var vs $N$ per initialisation width; check whether the exponential base moves toward 1 (plateau vanishes) or stays near $3/8$.

\item \textbf{Noise-induced barren plateaus on top of the ansatz-induced one.} Wang et al.\ (2021) showed noise-induced barren plateaus for generic deep circuits, but the paper's analysis is noise-free (pure-state). Tensor-network-inspired ansatze are proposed for near-term hardware, so noise resilience of the C1+C2 decay is a first-order practical question. \emph{Next steps:} repeat the MC estimator on PennyLane's \texttt{default.mixed} device with a depolarising channel of strength $p \in \{0.001,\,0.01,\,0.05\}$ inserted after each 2-qubit block; fit Var vs $N$ per noise level; check whether the exponential base changes and whether a noise-only regime (small $N$, large $p$) yields additional exponential decay.

\item \textbf{Connection to overparameterisation / lazy-training regime.} The paper analyses a single-parameter derivative in the moderately-parameterised regime. Overparameterisation (parameter count exceeding $2^{2N}$) has been shown to lift trainability landscapes into a lazy-training / NTK regime with modified Var-of-gradient scaling. Whether the qMPS plateau survives, softens, or transitions across this threshold is open. \emph{Next steps:} extend the qMPS staircase to deeper variants (repeated blocks, multi-layer bricks) and vary total parameter count $M$ at fixed $N$; compute Var$[\partial]$ as a function of $M/2^{2N}$; check for a phase transition in the Var-vs-$N$ scaling as one crosses the overparameterisation threshold; compare to our fixed-depth results.

\end{enumerate}
